\documentclass[11pt]{article}
\usepackage[margin=27mm]{geometry}
\usepackage{amsmath,amssymb,amsthm,booktabs,hyperref}
\newtheorem{lemma}{Lemma}\newtheorem{theorem}[lemma]{Theorem}
\newcommand{\PSL}{\operatorname{PSL}}\newcommand{\PGL}{\operatorname{PGL}}
\newcommand{\N}{\operatorname{N}}\newcommand{\F}{\mathbb F}
\title{Subfield parabolics and affine overgroups}
\author{Human Compression Project}\date{}
\begin{document}\maketitle
\begin{abstract}
The stabilizer of a subfield line determines its projective overgroup,
while its affine overgroups are controlled by invariant vector spaces.
In even characteristic, one-dimensional quotients supply the additional
geometries omitted by the source. For an odd extension in odd
characteristic the available scalar group is too small to do so.
The quadratic projective case has exactly two conjugacy classes of
affine configurations, interchanged by an outer automorphism. We give
these arguments in common coordinates and then treat the small
exceptions separately.
\end{abstract}

\section{Premises and the representation interface}
Let $K=\F_s\subset F=\F_q$ and let $G=\PSL(2,F)$ act on
$F\cup\{\infty\}$. For an additive subgroup $V\le F$, write
$U_V=\{x\mapsto x+v:v\in V\}$. Let $C_D$ denote multiplication
by elements of a subgroup $D\le F^*$. The scalars present in the
stabilizer of $\infty$ in $G$ form
\[
 A_F=\begin{cases}F^*,&\operatorname{char}F=2,\\
 (F^*)^2,&\operatorname{char}F\text{ odd}.
 \end{cases}
\]
All products $U_V:C_D$ below use these literal transformations.

We retain the classical Dickson subgroup classification, its subfield
embedding and conjugacy statements, and its small exceptional group
normalizers. In particular, odd-degree $\PSL(2,K)$ subgroups form
one conjugacy class, as do the even-characteristic subfield groups.
The two classes of $\PGL(2,K)$ in $\PSL(2,\F_{s^2})$ are interchanged
by $\PGL(2,\F_{s^2})$. We do not infer configuration counts from these
subgroup counts alone.

The following is the explicit interface with the separate common
representation proof. It is not an import of the source's unpublished
Proposition 4.6: for a subgroup of projective type, a two-transitive
coset action has its natural Borel stabilizer, or one of the following
exceptional stabilizers:
\[
\begin{array}{c|c}
\text{group}&\text{additional point stabilizer}\\\hline
\PSL(2,4)\cong\PSL(2,5)&D_{10}\text{ or }A_4\text{, as applicable}\\
\PSL(2,7)&S_4\\
\PSL(2,9),\ \PSL(2,11)&A_5\\
\PGL(2,t)\ (t\text{ odd})&\PSL(2,t)\text{ (degree two)}\\
\PGL(2,5)&S_4
\end{array}
\]
The familiar $\PSL(2,2)=S_3$, $\PSL(2,3)=A_4$ and
$\PGL(2,3)=S_4$ actions are included with their elementary kernels.
The same common proof supplies the small exceptional overgroup lists
when the point stabilizer is $A_4,S_4,A_5$ or a small dihedral group.
This interface is proved in the separately reviewed manuscript
\emph{Two-transitive coset actions in dimension two}.
The coordinate and uniqueness lemmas below are independent of it.

\section{A fixed Borel has a unique projective overgroup}
\begin{lemma}\label{normalizer}
Suppose $|K|=s>2$ and $D\le K^*$ is nontrivial. For
$B=U_K:C_D\le G$,
\[
 \N_G(B)=U_K:C_{K^*\cap A_F}.
\]
In particular $\N_G(B)=B$ in each of these cases:
\begin{enumerate}
\item $s\ge4$ is even, $D=K^*$;
\item $s\ge5$ is odd, $[F:K]$ is odd, $D=(K^*)^2$;
\item $s\ge3$ is odd, $[F:K]=2$, $D=K^*$.
\end{enumerate}
\end{lemma}
\begin{proof}
The unique Sylow subgroup $U_K$ for the characteristic of $F$ is
characteristic in $B$. It has the unique common fixed point $\infty$,
so a normalizer of $B$ has the form $x\mapsto ax+b$ with
$a\in A_F$. Normalizing $U_K$ gives $aK=K$, hence $a\in K^*$.
Conjugating a scalar $c\ne1$ in $D$ gives translation term
$(1-c)b$. This belongs to $K$ exactly when $b\in K$. Conversely
these conditions plainly normalize $B$, proving the formula.
For an odd-degree extension, an element of $K^*$ is a square in
$F$ exactly when it is a square in $K$; for a quadratic extension
every element of $K^*$ is a square in $F$. The assertions follow.
\end{proof}

\begin{lemma}\label{unique}
In the three cases of Lemma~\ref{normalizer}, respectively put
$M=\PSL(2,K)$, $M=\PSL(2,K)$ and $M=\PGL(2,K)$.
Assume the extension degree is prime, so that the specified $M$ is
maximal in $G$. Then exactly one conjugate subfield group of this
type contains the fixed $B$ as a natural Borel.
\end{lemma}
\begin{proof}
The Borel subgroups inside $M$ form one $M$-class and are
self-normalizing in $M$. For any ambient conjugacy class of copies
of $M$ that contains $B$, their set above $B$ is a transitive
$\N_G(B)$-set: adjust a transporter inside the target $M$ using
its internal Borel conjugacy. Its stabilizer is $B$, because $M$
is maximal and self-normalizing. Lemma~\ref{normalizer} makes this
set a singleton.

In the first two cases there is only one ambient class of $M$.
In the third case there are two. The Borel subgroups have two ambient
classes as well, which can be seen directly. After moving their
unique fixed point to $\infty$, their radicals are the subfield
lines $aK$. Conjugation multiplies $a$ by an element of $A_F$.
As $K^*\subset A_F$, the quotient $F^*/(A_F K^*)$ has order two.
The two classes are interchanged by any nonsquare scalar in
$\PGL(2,F)$. That same outer coset interchanges the two classes
of $M$. Since one class of $M$ contains only one class of Borels,
the two classes of $M$ must contain different Borel classes. Thus
only one class of $M$ can contain the fixed $B$, completing the proof.
\end{proof}

Here the radical description is literal: its scalar complement
$K^*$ forces an additive radical of order $s$ to be a one-dimensional
$K$-space. Any translated scalar complement can be made linear by
an affine translation. The same argument gives one Borel class in
the first two cases: $A_F K^*=F^*$ for an odd-degree odd-characteristic
extension, and $A_F=F^*$ in even characteristic.

\section{All affine covers of the fixed Borel}
\begin{lemma}\label{affine}
In any of the three cases above, let $H$ be an affine subgroup of
$G$ containing $B$ properly, with $B$ maximal in $H$.
If $D=K^*$, then the two-transitive possibilities are exactly
\[
 H=U_V:C_{K^*},\qquad K\subset V\subset F,
 \qquad \dim_K V=2.
\]
If $D=(K^*)^2$ and $s\ge5$ is odd in an odd-degree extension,
there is no such two-transitive affine overgroup.
\end{lemma}
\begin{proof}
An affine group containing $U_K$ must fix $\infty$, since that
is its only possible common fixed point. Write $U_V$ for the
translation kernel of $H$. The group $D$ generates the field $K$:
this is immediate for $K^*$; for $(K^*)^2$ its order is
$(s-1)/2>\sqrt{s}-1$, larger than the multiplicative group of any
proper subfield. Thus $V$ is a $K$-space containing $K$.

We first see that the scalar part of $H$ can be taken literally.
If $x\mapsto ax+b$ belongs to $H$, commute it with a fixed scalar
$c\ne1$ in $D$. The resulting translation is a nonzero $K$-multiple
of $b$, so $b\in V$. Removing this translation gives the scalar
$a$ itself in $H$. Therefore $H=U_V:C_{D'}$ for some
$D\le D'\le A_F$ stabilizing $V$.
If $V=K$, then $D'\le K^*\cap A_F=D$, contrary to properness.
If $V>K$, the intermediate subgroup $U_V:C_D$ properly contains
$B$, so maximality forces $D'=D$.

Identify the cosets of $B$ in $H$ with the additive quotient $V/K$.
Translations act regularly and the point stabilizer induces exactly
the scalar group $D$. Hence two-transitivity is equivalent to
transitivity of $D$ on the nonzero vectors of $V/K$. If $D=K^*$,
this holds exactly when $\dim_K(V/K)=1$: the orbit has size at most
$s-1$, and a one-dimensional quotient has exactly $s-1$ nonzero
vectors. For $D=(K^*)^2$, even a one-dimensional quotient has more
nonzero vectors than $|D|$, so transitivity is impossible.
\end{proof}

The scalar case with $D=K^*$ is the previously accepted affine-plane
construction. Its complete parametrization and Frobenius action are
reused from \emph{Scalar orbits and subfield configurations}; they
are not a new acquisition here. Lemma~\ref{affine} explains why no
other affine radical or enlarged complement can occur.

\begin{theorem}\label{borelconfigs}
Assume the representation interface in Section 1.
\begin{enumerate}
\item For even $s\ge4$ and prime $[F:K]=m$, the configurations
with $M=\PSL(2,K)$ and natural Borel intersection are precisely
$\{M,U_V:C_{K^*}\}$ with $\dim_K V=2$ and $K\subset V$.
They have $(s^{m-1}-1)/(s-1)$ conjugacy classes. Their full
semilinear classes are the Frobenius orbits on the same projective
parameter set. At $m=2$ both counts are one.
\item For odd $s\ge5$ and odd prime extension degree, there is no
configuration with $M=\PSL(2,K)$ and its natural Borel intersection.
\item For odd $s\ge3$ and $F=\F_{s^2}$, there are exactly two conjugacy
classes and one automorphism class of configurations with
$M=\PGL(2,K)$ and natural Borel intersection. Their other parabolic
is $U_F:C_{K^*}$.
\end{enumerate}
\end{theorem}
\begin{proof}
The representation interface and the characteristic restrict a
non-affine overgroup of the given Borel to the same projective type
as $M$. Indeed, the natural Borel stabilizer determines its
field from its characteristic part and scalar order. A cyclic group
cannot contain $B$. A dihedral candidate can occur only for the
small coincidences $B=D_{10}$ in characteristic five or $B=D_6$
in characteristic three; a proper dihedral overgroup would require
a cyclic element with mixed characteristic order, or order nine in
characteristic three, neither of which occurs in a projective group
in that characteristic. The exceptional groups $A_4,S_4,A_5$
are checked by their orders and the small-overgroup interface.
The small coincidence $U_4:C_3=A_4$ gives $A_5=\PSL(2,4)$
in characteristic two; the coincidence $U_5:C_2=D_{10}$ gives
$A_5=\PSL(2,5)$ in characteristic five. They introduce no new
projective overgroup. Lemma~\ref{unique} excludes a second such
copy. The remaining possibilities are exactly those in
Lemma~\ref{affine}.

For the first case, the accepted scalar-orbit proof gives the stated
counts and parameter set. In the second case there is no affine
possibility. In the third case the only two-dimensional $K$-space
is $F$ itself. Each fixed Borel therefore has exactly one projective
and one affine overgroup, so its two ambient classes give exactly
two configuration classes. A nonsquare scalar exchanges the two,
proving the single automorphism class. All intersections are exactly
$B$; maximality of $M$ gives generation of $G$, and the affine action
on $F/K$ supplies local two-transitivity.
\end{proof}

\section{Odd subfields and the exceptional stabilizers}
Here a subgroup defined over $K$ means that every projective element
lies in $\PGL(2,K)$ and is individually fixed by the Frobenius
of $F/K$. Mere stability of the subgroup as a set is insufficient.

\begin{lemma}\label{descent}
Let $K\subset F$ be finite fields.
Suppose $L_1,L_2\le\PGL(2,K)$ have trivial centralizer in
$\PGL(2,F)$. Every element of $\PGL(2,F)$ conjugating $L_1$
to $L_2$ belongs to $\PGL(2,K)$. If the extension degree is odd, two such subgroups
of $\PSL(2,K)$ that are not conjugate there cannot become conjugate
in $\PSL(2,F)$.
\end{lemma}
\begin{proof}
Let $gL_1g^{-1}=L_2$ and let $\sigma$ be the Frobenius generating
$\operatorname{Gal}(F/K)$. For every $h\in L_1$, both $h$ and
$ghg^{-1}$ are fixed by $\sigma$. Thus $g^{-1}\sigma(g)$
centralizes $L_1$ and is the identity. A projective matrix fixed by
$\sigma$ has a representative over $K$: normalize one nonzero
entry to one, after which all entry ratios are fixed. Hence
$g\in\PGL(2,K)$. Finally
$\PGL(2,K)\cap\PSL(2,F)=\PSL(2,K)$ for odd extension degree,
by the square criterion used in Lemma~\ref{normalizer}.
\end{proof}

\begin{lemma}\label{unique-descent}
Let $M$ be the standard $\PSL(2,K)$ or $\PGL(2,K)$ in
$\PGL(2,F)$. A subgroup $B\le M$ with trivial projective
centralizer is contained in no other $\PGL(2,F)$-conjugate of $M$.
\end{lemma}
\begin{proof}
Suppose $B\le gMg^{-1}$. Then $g^{-1}Bg$ and $B$ are both
subgroups of $\PGL(2,K)$, and their centralizers are trivial.
Lemma~\ref{descent} gives $g\in\PGL(2,K)$. That group normalizes
both possible choices of $M$, so $gMg^{-1}=M$.
\end{proof}

In the odd-characteristic fields used in this section, the groups
$A_4,S_4,A_5$ have trivial projective centralizer.
Here is an elementary check. In
characteristic three choose two distinct cyclic subgroups of order
three: their unipotent centralizers have different fixed points and
intersect trivially. Indeed, two distinct such subgroups with the
same fixed point would commute and generate a group of order nine,
which cannot lie in any of the three groups under consideration. In other characteristics the centralizer of each
such subgroup is its torus. An element centralizing both must lie
in both tori. A nonidentity noninvolution has its own unique torus
centralizer; an involution has a dihedral centralizer whose odd-order
subgroups lie in its torus. In either case both subgroups of order
three would be the same unique subgroup of order three in that
torus, a contradiction.

\begin{theorem}\label{odd}
Assume the common representation interface. For $q=s^m$ odd with
$m$ an odd prime and maximal $M=\PSL(2,s)$, the only configurations
are those with $s=3$ and two $A_4$ parabolics meeting in $C_3$.
Their counts are the corrected counts in the accepted scalar-orbit
proof. In particular no additional configurations arise from the
exceptional $A_4,S_4,A_5$ stabilizers at $s=5,7,9,11$.
\end{theorem}
\begin{proof}
For $s\ge5$ the natural Borel case is excluded by
Theorem~\ref{borelconfigs}. The exceptional point stabilizer is
$A_4$ at $s=5$, $S_4$ at $s=7$, and $A_5$ at $s=9,11$.
The small overgroup interface in the actual characteristic leaves
only another copy of the same $M$ as a two-transitive overgroup.
All copies of this subfield group are conjugate under
$\PGL(2,F)$, by the retained subfield classification. The exceptional
stabilizers have trivial projective centralizer, as just proved.
Lemma~\ref{unique-descent} therefore gives a unique $M$ containing
the fixed stabilizer, ruling out a second parabolic. This argument
also explains why simply fusing internal subgroup classes would not
be an adequate replacement for a literal overgroup proof.

At $s=3$, $B=C_3$ is unipotent. Its normalizer in $G$ is $U_F$:
an affine multiplier preserving the translation line $\F_3$ is
$\pm1$, and $-1$ is a nonsquare in $F$ because $m$ is odd.
The elementary degree argument for a subgroup with point stabilizer
$C_3$ gives $C_6,S_3,A_4$. The first two would normalize $C_3$
and contain an involution, impossible in $U_F$. Only the $A_4$
overgroups remain. Their complete conjugacy and Frobenius-orbit
counts, with unordered interchange, were proved in the accepted
scalar-orbit component.
\end{proof}

\section{The small even-field cases}
\begin{theorem}\label{small-even}
The case $M=\PSL(2,2)$ in $G=\PSL(2,4)$ gives exactly one
conjugacy and one automorphism class of each of the following types:
\[
 (S_3,S_3;C_2),\qquad(S_3,C_2^2;C_2),\qquad(S_3,A_4;C_3).
\]
For $M=\PSL(2,4)$ in $G=\PSL(2,4^m)$ with $m$ prime,
the exceptional intersection $D_{10}$ gives the single type
$(A_5,D_{30};D_{10})$, with one class in both senses, exactly
when $15\mid4^m\pm1$. This condition is equivalent here to $m=2$.
\end{theorem}
\begin{proof}
Use the natural $A_5$ model for $\PSL(2,4)$. A subgroup $S_3$
is the normalizer of a three-cycle subgroup; its involutions are
products of two transpositions. It is maximal in $A_5$. There are
ten such subgroups, fifteen involutions, five point stabilizers
$A_4$, and five subgroups $C_2^2$; each latter group is the
centralizer of any of its nonidentity elements. These counts also
follow directly by choosing the supports of cycles in five letters.

A point stabilizer $C_2$ in a two-transitive action has degree two
or three, since $d-1\mid2$. Thus its overgroup has order four
or six. The order-four possibility in $A_5$ is $C_2^2$, uniquely
specified by the involution. The order-six possibility is $S_3$.
Each involution belongs to exactly two $S_3$ subgroups: count
$10\cdot3=30$ incidences and use transitivity on the fifteen
involutions. Their unordered pair is unique, proving the first
configuration count. For the different-type pair $S_3,C_2^2$,
all chains consisting of $S_3$ and one of its involutions are
conjugate, proving the second count.

For the intersection $C_3$, the small degree argument leaves
$C_6,S_3,A_4$. There is no $C_6$ in $A_5$, and the normalizer
$S_3$ of the fixed $C_3$ is its unique order-six overgroup.
Exactly two point stabilizers $A_4$ contain this $C_3$, corresponding
to its two fixed letters. Its $S_3$ normalizer interchanges these
two letters. Hence there is exactly one configuration of the third
type. Each asserted intersection and each local action follows from
the indicated elementary subgroup lattices. Maximality of the
$S_3$ parabolic gives generation and connectedness.

Now let $q=4^m$. The two-transitive overgroup interface for $D_{10}$
leaves $A_5,D_{20},D_{30}$. The cyclic torus containing $C_5$
has odd order $q\pm1$. The normalizer of the specified $D_{10}$
is $D_{10}$ itself: a torus element normalizing it has square in
$C_5$, and squaring is invertible on the odd-order quotient torus.
There is one class of $A_5$ subfield groups and one internal class
of $D_{10}$, so the transporter argument gives a unique $A_5$
above the fixed intersection. An overgroup $D_{20}$ is impossible,
since it would contain an element of order ten in characteristic two.

A $D_{30}$ overgroup exists precisely when the same torus contains
$C_{15}$, and is then uniquely $\langle C_{15},D_{10}\rangle$.
The $D_{10}$ subgroups form one ambient class, proving both counts.
Their local indices are six in $A_5$ and three in $D_{30}$;
the corresponding natural actions are two-transitive. At $m=2$
the condition holds because $4^2-1=15$. If $m$ is odd, five
divides $4^m+1$, while three divides $4^m-1$, so no torus order
is divisible by fifteen. This proves the last assertion.
\end{proof}

\section{The remaining quadratic projective intersections}
\begin{theorem}\label{quadratic}
Let $s$ be odd and $G=\PSL(2,s^2)$, with maximal parabolic
$M=\PGL(2,s)$. An intersection $\PSL(2,s)$ gives no second
parabolic unless $s=3$, in which case the other parabolic is $A_5$.
That exceptional configuration has two conjugacy classes and one
automorphism class. The exceptional $S_4$ intersection at $s=5$
gives no configuration with a second parabolic.
\end{theorem}
\begin{proof}
Put $L=\PSL(2,s)$. Its ambient normalizer is $M$, by the
classical subfield normalizer formula. Any other $\PGL(2,s)$
overgroup would have $L$ as its characteristic subgroup of index
two, and would therefore lie in $\N_G(L)=M$. Hence it equals $M$.
The representation and small-overgroup interface leaves only $A_5$
as an additional possibility, and only when $s=3$, so that $L=A_4$.
For clarity, an overgroup of larger projective field cannot create
a hidden option: the subfield classification forces its field to
contain $\F_s$. Inside $\F_{s^2}$ the only larger possibility is
the entire ambient group. Its index over $L$ is $s(s^2+1)$,
which exceeds $|L|+1$ and violates the two-transitivity order bound.

For $s=3$, the group $G=A_6$ has two classes of $A_4$, with
normalizer $S_4$, and two classes of $A_5$. Each $A_5$ has one
internal class of $A_4$. An outer automorphism interchanges the two
classes of each type; hence the two $A_5$ classes contain different
$A_4$ classes. Fix $L=A_4$. The members of its one applicable
$A_5$ class containing it form a transitive $\N_G(L)=S_4$-set
with stabilizer $L$: the $A_5$ is self-normalizing and its
intersection with $S_4$ is precisely $L$. Thus there are two
such $A_5$ overgroups, exchanged by the fixed $S_4$ parabolic.
This gives one configuration class for each of the two $A_4$
classes, and therefore two in all. The outer automorphism fuses them,
giving one automorphism class. The local indices are two and five,
with the standard two-transitive actions.

At $s=5$, the subgroup $S_4$ is the exceptional point stabilizer
in $M\cong S_5$. The common overgroup interface leaves only
$\PGL(2,5)$. Its $S_4$ stabilizer has trivial projective centralizer and
is defined over $\F_5$. Lemma~\ref{unique-descent} gives a unique
$\PGL(2,5)$ overgroup, ruling out a second parabolic.
\end{proof}

The $q=9$ configuration with two $S_4$ parabolics and intersection
$D_8$ is imported from the separately accepted duad--syntheme proof,
which also identifies its full graph group. It is not proved again here.

\section{Source claims, corrections and attribution}
The original numerical rows are interpreted under their stated
maximal-parabolic hypotheses. The results above give the following
joins; the common representation interface remains a separately
reviewed dependency.
\begin{itemize}
\item L063 Lemma 5.11 and Proposition 5.12, and L064 Lemma 21 /
Proposition 22: Theorem~\ref{borelconfigs} recovers Table 4 row 1
at $m=2$ and supplies the missing affine families for other prime
degrees. The published restriction to $m=2$ is false. The fixed
appendix's Lemma 21 already lists the affine type without that
restriction, so its statement is distinguished from the false
exhaustivity conclusion of Proposition 22.
\item Table 4 rows 2--4 are Theorem~\ref{small-even} at $q=4$.
Row 5 repeats Table 1 row 1 and is the $D_{10}$ part of that theorem.
The natural-Borel subcase at $s=4$ also reuses the accepted Table 1
rows 2--4 scalar-orbit counts.
\item L063 Lemmas 5.13--5.15 / Proposition 5.16, and L064
Lemmas 23--25 / Proposition 26: Theorem~\ref{odd} proves the
remaining exclusion arguments in their Section 5.4.2 context of
odd prime extension degree. The two surviving Table 5 rows retain
the accepted correction for unordered inner counts. The statement
of Lemma 5.13 is read with that surrounding odd-degree hypothesis.
\item L063 Lemma 5.18 / Proposition 5.21 and L064 Lemma 28 /
Proposition 30: Table 6 row 1 is Theorem~\ref{borelconfigs}(3).
Theorem~\ref{quadratic} gives row 2 and the nonexistence conclusions
above $\PSL(2,s)$ and, at $q=25$, above $S_4$. These also cover
L063 Lemmas 5.19--5.20 and L064 Lemma 29. The fixed appendix's
Proposition 30 does not include a separate $q=25$/$S_4$ subcase;
that published-source obligation is proved explicitly here.

The proof of Lemma 24 in the fixed appendix also states that
$\PSL(2,3^{2m})$ contains no $\PGL(2,3^w)$ subgroup.
That auxiliary assertion is false: $w=1$ and $w=m$ are allowed
subfield embeddings. It is not used here. In the latter case,
$5$ does not divide $|\PGL(2,3^m)|$ for odd $m$, so it cannot
contain $A_5$; the full exclusion is supplied by the explicit
representation interface.
\item Table 6 row 3, L063 Lemma 5.17 and L064 Lemma 27 are joined
to the accepted $q=9$ component, together with the small overgroup
list. Their full graph properties are already covered separately.
\end{itemize}

The original family constructions and classical subgroup theory retain
their attribution. The coordinate affine-cover argument and the
Galois-transporter descent argument are the common proofs used here;
the earlier scalar-orbit and exceptional-geometry manuscripts are
explicitly reused. This manuscript does not claim a new principal
refutation, a new answer to an open conjecture, Lean verification or
a compression measurement. Whole-original disposition requires the
remaining source components and their exact claim ledgers.

\begin{thebibliography}{9}
\bibitem{source} F.~Buekenhout, J.~De Saedeleer and D.~Leemans,
\emph{On the rank two geometries of the groups $\PSL(2,q)$: part II},
Ars Mathematica Contemporanea \textbf{6} (2013), 365--388,
Sections 5.4--5.5 and Tables 4--6.
\bibitem{appendix} The same authors, \emph{Part II (Appendix)},
retained 21-page edition, Lemmas 21--29 and Propositions 22, 26, 30.
\end{thebibliography}
\end{document}
